\chapter{Introduction}
\label{chap:introduction}

\section{Background}
\label{sec:background}

The Large Hadron Collider (LHC) at CERN collides protons at a
centre-of-mass energy of up to $14$\,TeV and records the debris with
detectors such as ATLAS and CMS. Most of the interesting particles produced
in these collisions are not observed directly; instead they decay into
sprays of hadrons that are reconstructed into collimated bundles called
\emph{jets}~\citep{cacciari2008,salam2010}. Reading the internal structure
of a jet --- its \emph{substructure} --- is the key to deciding which kind
of particle produced it~\citep{larkoski2020,kogler2019}.

The top quark is the heaviest known elementary particle, with a mass of
$m_t\simeq172.76$\,GeV~\citep{pdg2022}. It decays almost exclusively through
$t\rightarrow Wb$, and when the $W$ boson decays hadronically the full chain
becomes $t\rightarrow Wb\rightarrow q\bar{q}'b$. If the top quark is produced
with a large transverse momentum (it is then said to be \emph{boosted}), the
three resulting quarks are captured inside a single large-radius jet that
carries a characteristic three-pronged energy pattern and a mass near
$m_t$. Distinguishing such ``top jets'' from the vastly more common jets
initiated by light quarks and gluons (Quantum Chromodynamics, QCD, jets) is
known as \emph{top tagging}, and it is one of the canonical classification
problems of collider physics~\citep{kasieczka2019landscape}.

Over the last decade, top tagging has become a benchmark for machine
learning in physics. Early approaches engineered physically motivated
variables such as the jet mass, $N$-subjettiness~\citep{thaler2011} and
energy-correlation functions~\citep{larkoski2013} and fed them to boosted
decision trees or shallow networks. The field then moved to deep models that
act directly on the low-level constituents: convolutional networks on
pixelated ``jet images''~\citep{deoliveira2016,komiske2017}, graph networks
such as ParticleNet~\citep{qu2020particlenet}, and most recently
attention-based models such as the Particle
Transformer~\citep{qu2022particletransformer}. On the standard reference
dataset~\citep{kasieczka2019dataset}, these models reach areas under the ROC
curve (AUC) above $0.98$, and a community
comparison~\citep{kasieczka2019landscape} established the relative ordering
that is now treated as the state of the art.

\section{Problem Statement}
\label{sec:problem_statement}

The published top-tagging literature is almost entirely concerned with the
\emph{best achievable} discrimination, quoted at the full training-set size
of $1.2\times10^{6}$ jets and on large GPU clusters. Two questions of direct
practical importance are rarely addressed together:

\begin{enumerate}
    \item \textbf{Data efficiency.} Generating labelled Monte Carlo training
    samples is computationally expensive. How much does each architecture's
    performance depend on the amount of training data, and does the ranking
    of the models change in the small-data regime that a student or a small
    group actually works in?
    \item \textbf{Probability calibration.} A high AUC only certifies that a
    model \emph{ranks} signal above background. Many downstream physics
    analyses instead use the classifier score as a \emph{probability} (for
    example, in a likelihood fit). Are the scores produced by these models
    trustworthy probabilities, and if not, can they be cheaply corrected?
\end{enumerate}

This thesis addresses both questions on the public reference dataset using a
single, reproducible pipeline that can be run end-to-end on free hardware.

\section{Aim and Research Objectives}
\label{sec:objectives}

The overall \textbf{aim} of this project is to build a transparent,
reproducible top-tagging pipeline and use it to study the data efficiency
and probability calibration of four representative classifiers, rather than
to chase a new record AUC.

The specific \textbf{objectives} are:
\begin{enumerate}
    \item \textbf{To implement} a unified pipeline that trains and evaluates
    four architectures --- a gradient-boosted decision tree (BDT) on
    engineered features, a convolutional neural network (CNN) on jet images,
    the ParticleNet graph network, and the Particle Transformer --- under
    identical preprocessing, data splits and random seeds.
    \item \textbf{To reproduce} the established performance ordering on the
    Top-Quark Tagging Reference Dataset and report AUC, accuracy and
    background rejection at fixed signal efficiency.
    \item \textbf{To quantify data efficiency} by re-training each model from
    scratch at four training-set sizes spanning a decade
    ($5\times10^{3}$ to $5\times10^{4}$ jets) and measuring how performance
    scales.
    \item \textbf{To assess probability calibration} using reliability
    diagrams and the Expected Calibration Error (ECE), and to test whether a
    single-parameter temperature rescaling restores calibration.
    \item \textbf{To make the work fully reproducible}, releasing the code,
    trained models and documents as open source and providing a cloud-GPU
    workflow that any student can follow.
\end{enumerate}

\section{Research Questions}
\label{sec:research_questions}

\begin{enumerate}
    \item Does the literature ordering (Particle Transformer $>$ ParticleNet
    $>$ CNN $>$ BDT) hold when all four models are trained identically on a
    modest data and compute budget?
    \item How does each model's discrimination scale with the training-set
    size, and is there a regime in which the simple BDT is preferable to the
    deep models?
    \item Are the classifier scores well-calibrated probabilities, and does
    temperature scaling fix any miscalibration?
\end{enumerate}

\section{Significance of the Study}
\label{sec:significance}

This study is significant for three reasons. First, it provides a clear,
honest, reproducible reference implementation of four top-tagging
architectures that can be used for teaching and as a baseline for future
projects at the host institution and beyond. Second, the data-efficiency
results give practitioners concrete guidance on which model to choose under a
limited Monte Carlo budget --- a situation common outside the large LHC
collaborations. Third, the calibration analysis highlights a subtle but
important point often overlooked in physics machine-learning papers: a model
with the best ranking power is not necessarily the one whose scores can be
read as probabilities. Demonstrating that the entire study can be carried out
on a free cloud GPU also lowers the barrier to entry for students in
resource-constrained settings.

\section{Scope and Limitations}
\label{sec:scope}

The study is deliberately bounded. We use stratified subsamples of up to
$5\times10^{4}$ jets per split rather than the full $1.2\times10^{6}$-jet
dataset, so the absolute AUC values are a little below the published
saturation values; the \emph{trends} and \emph{rankings}, not the absolute
records, are the object of study. The dataset itself is simulated with a
single detector parametrisation and contains no pile-up, so questions of
detector robustness and high-luminosity performance are out of scope. The
four architectures are representative rather than exhaustive, and each is
trained with a fixed, modest hyper-parameter configuration rather than an
exhaustive search. These limitations are revisited in
Chapter~\ref{chap:conclusions}.

\section{Thesis Organization}
\label{sec:organization}

The remainder of this thesis is organised as follows.
Chapter~\ref{chap:literature_review} reviews the physics of jets and top
quarks and the machine-learning approaches to tagging.
Chapter~\ref{chap:methodology} describes the dataset, the feature
engineering, the four model architectures and the experimental protocol.
Chapter~\ref{chap:results} presents the headline benchmark, the
data-efficiency sweep and the calibration analysis, and discusses their
meaning. Chapter~\ref{chap:conclusions} summarises the findings, states the
limitations and proposes future work. Appendix~\ref{chap:appendix_a} gives
reproduction instructions and selected implementation details.
